\chapter{Save this session as a script (menu 8)}
\label{ch:menu8}
\menulabel{8}

\section{Purpose}

Keep every input you have typed so far as a replayable script: record once,
replay forever.

\section{What you need}

Nothing but the session itself. The alternative route is to start the program
with \cmd{--record FILE}, which records from the first prompt.

\section{Walkthrough}

\lstinputlisting[style=fbkcode, firstline=174, lastline=181,
  title={Saving a session at menu 8 (from the recording example)}]
  {examples/expected/m8_script.txt}

The saved file contains the six input lines typed \emph{before} the saving
action -- the save-path answer itself is not part of it:

\lstinputlisting[style=fbkcode,
  title={The recorded script (\file{examples/expected/saved\_session.txt})}]
  {examples/expected/saved_session.txt}

\section{What you get}

A text file with one input per line. Its format is the script format of
Chapter~\ref{ch:running}: an empty line means ``press Enter'' (accept the
default), \code{\#} starts a comment, and everything else goes to the prompt.

\section{How to read the numbers}

\begin{readbox}
The count in ``Saved $N$ input line(s)'' is the size of the recorded prefix.
When you replay a script, the program prints the same conversation again --
prompts plus your recorded answers -- so the replay file is a complete,
diffable transcript. Replaying the same script twice is byte-identical; that
property is a tested guarantee, not a hope.
\end{readbox}

\section{Worked example}

A recording ends at the moment of saving, so replaying it runs the recorded
analyses and then finishes the input at the save prompt:

\lstinputlisting[style=fbkcode, firstline=178, lastline=182,
  title={Replaying the saved script: the recorded part runs, then the input ends}]
  {examples/expected/m8_replay.txt}

To record a complete workflow, either put \menu{8} last (the saved prefix is
then everything that matters), or use \cmd{--record FILE}, which appends the
whole session including the exit.

\section{Caveats, refusals and related sections}

\begin{refusalbox}
\begin{itemize}
  \item A script replays inputs, not files: if a recorded script points at
    \file{work/job.out}, that file must exist when you replay it.
  \item Replaying is exactly equivalent to typing; it grants no extra
    capabilities (nothing runs calculations, nothing touches the network).
  \item If the input runs out mid-session the program says
    ``Cancelled (input finished)'' and continues cleanly -- it never blocks
    waiting for a terminal that is not there.
\end{itemize}
\end{refusalbox}

Related: Chapter~\ref{ch:running} (the full script model and the CLI), and
\file{docs/manual/examples/} in the repository, whose scripts are replayed by
the test suite.
